%****************************************************
%	CHAPTER 1 - INTRODUCTION
%****************************************************
\chapter{Introduction}
\label{ch:ch1.introduction}
%====================================================
\section{Background}
\label{sec:ch1.background}
%----------------------------------------------------

*need to add section discussing pancake ice environmental factors and why were looking for what were looking at*

%====================================================
\section{Problem Statement}
\label{sec:ch1.problemstatement}
%----------------------------------------------------

The \ac{MIZ} is a highly dynamic and complex environment, where the interaction between ocean waves and sea ice leads to a variety of ice types and formations. Knowledge the characteristics of sea ice in the \ac{MIZ} is crucial for improving our understanding of not only the processes that govern the \ac{MIZ}, but also the broader environmental systems that it influences. In-situ observations of the \ac{MIZ} are challenging due to many factors, including harsh weather conditions and its remote location. These challenges have led to certain regions and time periods being under-sampled, resulting in biases and gaps in our understanding of the \ac{MIZ} as a whole. \ac{LiDAR} technology has become more commonplace in Earth observation applications in recent years, but has not been widely used for sea ice observations in the \ac{MIZ}. It has potential to provide three-dimensional data on sea ice characteristics, but the feasibility of using \ac{LiDAR} for this application must be further investigated. 

%====================================================
\section{Project Objectives}
\label{sec:ch1.projectobjectives}
%----------------------------------------------------

This research project aims to investigate the feasibility of using \ac{LiDAR}, through the use of the Livox Avia sensor, for ship-based observation of sea ice in the \ac{MIZ} and the extraction of useful characteristics from the resulting point cloud data. The objectives of this project are as follows:
\begin{itemize}
    \item Identify what characteristics of sea ice in the \ac{MIZ} can be observed through the use of a \ac{LiDAR} sensor.
    \item Identify potential methods for the extraction of these characteristics from \ac{LiDAR} point cloud data.
    \item Characterise the Livox Avia sensor to determine its strengths and limitations in the context of observing sea ice in the \ac{MIZ} from a research vessel.
    \item Develop and evaluate a data processing pipeline for the extraction of sea ice characteristics from point cloud data.
\end{itemize}

%====================================================
\section{Scope, Limitations and Assumptions}
\label{sec:ch1.scope}
%----------------------------------------------------

The scope of this project is limited to the processing and analysis of \ac{LiDAR} point cloud data collected using the Livox Avia sensor from the Aalto University Wave and Ice Tank, during the \ac{SCALE} Winter 2022 cruise, and data collected during testing and characterisation of the sensor under controlled conditions.

No field deployments of the Livox Avia sensor were conducted as part of this project, with all in-situ data collection conducted by other parties prior to this investigation. As such, any limitations or issues present in the collected data are outside the control of this project, as well as any limitations imposed by the deployment platform or sensor configuration used during data collection. 

MATLAB was selected as the programming environment for this project due to prior experience and its extensive suite of toolboxes for point cloud processing, data visualisation and analysis. The use of MATLAB may limit the portability of developed code to other programming environments, and may also limit performance when processing large datasets compared to lower-level programming languages.

%====================================================
\section{Plan of Development}
\label{sec:ch1.planofdevelopment}
%----------------------------------------------------


%----------------------------------------------------
%****************************************************
% END
%****************************************************
